% ==============================================================================
% EDC BOOK 2: WEAK SECTOR — CANONICAL SPINE
% ==============================================================================
% This is the ONLY reader-facing path through Book 2.
% All content flows through this file in coherent narrative order.
%
% Structure:
%   Part I   (Ch 1-6):  Physical Picture — Intuition and case studies
%   Part II  (Ch 7-12): Predictions — Electroweak observables and mixing
%   Part III (Ch 13-16): Machinery — OPR derivation chain
%   Epilogue (Ch 17):   Nuclear Applications Preview
%
% Created: 2026-01-29
% ==============================================================================

\documentclass[11pt,a4paper]{book}

% ═══════════════════════════════════════════════════════════════════════════════
% PACKAGES (inherited from rebuild)
% ═══════════════════════════════════════════════════════════════════════════════
\usepackage{fontspec}
\usepackage{amsmath,amssymb,amsthm}
\usepackage{physics}
\usepackage{geometry}
\usepackage{graphicx}
\usepackage{float}
\usepackage[dvipsnames]{xcolor}
\usepackage{tcolorbox}
\usepackage{hyperref}
\usepackage[backend=biber,style=numeric-comp,sorting=none,maxbibnames=3]{biblatex}
\addbibresource{bib/part2_backbone.bib}
\usepackage{tikz}
\usetikzlibrary{calc,angles,quotes,decorations.markings,decorations.pathmorphing,positioning}
\usetikzlibrary{shapes,arrows,arrows.meta,shapes.geometric,fit,backgrounds}
\usepackage{multirow}
\usepackage{booktabs}
\usepackage{longtable}
\usepackage{enumitem}
\usepackage{fancyhdr}
\usepackage{titlesec}
\usepackage{epigraph}
\usepackage{array}
\usepackage{colortbl}
\usepackage{pifont}
\usepackage[normalem]{ulem}
\usepackage{tabularx}
\usepackage{adjustbox}

\tcbuselibrary{breakable,skins}

% ═══════════════════════════════════════════════════════════════════════════════
% PAGE SETUP
% ═══════════════════════════════════════════════════════════════════════════════
\geometry{margin=1in}
\raggedbottom

\pagestyle{fancy}
\fancyhf{}
\fancyfoot[C]{\thepage}
\renewcommand{\headrulewidth}{0pt}

% ═══════════════════════════════════════════════════════════════════════════════
% THEOREM ENVIRONMENTS
% ═══════════════════════════════════════════════════════════════════════════════
\newtheorem{theorem}{Theorem}[chapter]
\newtheorem{lemma}[theorem]{Lemma}
\newtheorem{proposition}[theorem]{Proposition}
\newtheorem{corollary}[theorem]{Corollary}
\newtheorem{conjecture}[theorem]{Conjecture}
\theoremstyle{definition}
\newtheorem{definition}[theorem]{Definition}
\newtheorem{example}[theorem]{Example}
\newtheorem{postulate}[theorem]{Postulate}
\theoremstyle{remark}
\newtheorem{remark}[theorem]{Remark}
\newtheorem{observation}[theorem]{Observation}
\newtheorem{axiom}[theorem]{Axiom}

% ═══════════════════════════════════════════════════════════════════════════════
% CUSTOM COMMANDS
% ═══════════════════════════════════════════════════════════════════════════════
\newcommand{\Membrane}{\Sigma}
\newcommand{\Bulk}{\mathcal{M}_5}
\newcommand{\vscan}{v_{\text{scan}}}
\newcommand{\Rxi}{R_\xi}
\newcommand{\EDCPAPERS}{../../edc_papers}

% ═══════════════════════════════════════════════════════════════════════════════
% EPISTEMIC TAGS
% ═══════════════════════════════════════════════════════════════════════════════
\definecolor{derivedColor}{RGB}{0,100,0}
\definecolor{postulatedColor}{RGB}{150,0,0}
\definecolor{calibratedColor}{RGB}{0,0,150}

\newcommand{\tagBL}{\textcolor{brown}{\textbf{[BL]}}}
\newcommand{\tagDer}{\textcolor{derivedColor}{\textbf{[Der]}}}
\newcommand{\tagDc}{\textcolor{calibratedColor}{\textbf{[Dc]}}}
\newcommand{\tagI}{\textcolor{gray}{\textbf{[I]}}}
\newcommand{\tagCal}{\textcolor{postulatedColor}{\textbf{[Cal]}}}
\newcommand{\tagP}{\textcolor{postulatedColor}{\textbf{[P]}}}
\newcommand{\tagM}{\textcolor{gray}{\textbf{[M]}}}
\newcommand{\tagDef}{\textcolor{blue!60!black}{\textbf{[Def]}}}
\newcommand{\tagOPEN}{\textcolor{red!70!black}{\textbf{[OPEN]}}}
\newcommand{\tagOpen}{\textcolor{red!70!black}{\textbf{[OPEN]}}}
\newcommand{\tagNoGo}{\textcolor{red!70!black}{\textbf{[NO-GO]}}}

\newcommand{\degree}{°}

% ═══════════════════════════════════════════════════════════════════════════════
% TCOLORBOX STYLES
% ═══════════════════════════════════════════════════════════════════════════════
\tcbset{
  readerContract/.style={colback=blue!5, colframe=blue!50!black, fonttitle=\bfseries, title={Reader Contract}},
  guardrail/.style={colback=yellow!5, colframe=orange!60!black, fonttitle=\bfseries},
  falsifiability/.style={colback=red!5, colframe=red!40!black, fonttitle=\bfseries, title={Falsifiability Hooks}},
  mechanism/.style={colback=teal!5, colframe=teal!50!black, fonttitle=\bfseries}
}

\newtcolorbox{edcPostulateBox}[2]{colback=orange!5, colframe=orange!60!black, fonttitle=\bfseries, title={#1 \hfill {\small #2}}, breakable}
\newtcolorbox{edcDefinitionBox}[2]{colback=green!5, colframe=green!50!black, fonttitle=\bfseries, title={Definition: #1 \hfill {\small #2}}, breakable}
\newtcolorbox{edcPropositionBox}[2]{colback=purple!5, colframe=purple!50!black, fonttitle=\bfseries, title={Proposition: #1 \hfill {\small #2}}, breakable}
\newtcolorbox{edcLedgerBox}[2]{colback=blue!5, colframe=blue!50!black, fonttitle=\bfseries, title={Ledger: #1 \hfill {\small #2}}, breakable}
\newtcolorbox{edcConsequenceBox}[2]{colback=teal!5, colframe=teal!50!black, fonttitle=\bfseries, title={#1 \hfill {\small #2}}, breakable}
\newtcolorbox{edcWarningBox}[2]{colback=yellow!10, colframe=orange!70!black, fonttitle=\bfseries, title={Guardrail: #1}, breakable}
\newtcolorbox{edcLemmaBox}[2]{colback=cyan!5, colframe=cyan!50!black, fonttitle=\bfseries, title={Lemma: #1 \hfill {\small #2}}, breakable}
\newtcolorbox{edcTheoremBox}[2]{colback=blue!8, colframe=blue!60!black, fonttitle=\bfseries, title={Theorem: #1 \hfill {\small #2}}, breakable}
\newtcolorbox{edcCorollaryBox}[2]{colback=green!8, colframe=green!60!black, fonttitle=\bfseries, title={Corollary: #1 \hfill {\small #2}}, breakable}
\newtcolorbox{stepbox}[2]{colback=blue!5, colframe=blue!50!black, fonttitle=\bfseries, title={Step #1: #2}, breakable}
\newtcolorbox{questionbox}[1]{colback=yellow!10, colframe=orange!70!black, fonttitle=\bfseries, title={Question: #1}, breakable}
\newtcolorbox{answerbox}[1]{colback=green!5, colframe=green!50!black, fonttitle=\bfseries, title={Answer: #1}, breakable}
\newtcolorbox{gapbox}[1]{colback=red!5, colframe=red!50!black, fonttitle=\bfseries, title={Gap Identified: #1}, breakable}
\newtcolorbox{resolutionbox}[1]{colback=teal!5, colframe=teal!50!black, fonttitle=\bfseries, title={Resolution: #1}, breakable}

\newtcolorbox{edcAtAGlance}[1]{colback=blue!3, colframe=blue!50!black, fonttitle=\bfseries, title={At a Glance: #1}, before upper={\parskip=0.3em}, boxrule=1pt, arc=3pt}

\newcommand{\edcBaseline}[1]{\textbf{\textcolor{gray!70!black}{Standard Model \tagBL{}:}}\\#1\par\medskip}
\newcommand{\edcEDCView}[1]{\textbf{\textcolor{purple!70!black}{EDC Interpretation \tagP{}/\tagDc{}:}}\\#1\par\medskip}
\newcommand{\edcKeyInsight}[1]{\textbf{\textcolor{green!50!black}{Key Insight:}}\\\textit{#1}\par\medskip}
\newcommand{\edcFalsifiable}[1]{\textbf{\textcolor{red!60!black}{Falsifiable If:}}\\#1}

\tcbset{
    edcCornerstone/.style={colback=blue!5, colframe=blue!40!black, fonttitle=\bfseries},
    edcGuardrail/.style={colback=gray!5!white, colframe=gray!60!black, fonttitle=\bfseries},
    edcPPN/.style={colback=blue!5, colframe=blue!50!black, fonttitle=\bfseries},
    edcCanonical/.style={colback=yellow!5, colframe=orange!60!black, fonttitle=\bfseries},
    edcConcept/.style={colback=yellow!5, colframe=orange!50!black, fonttitle=\bfseries},
    edcPathway/.style={colback=purple!5, colframe=purple!40!black, fonttitle=\bfseries},
    edcModel/.style={colback=green!5, colframe=green!40!black, fonttitle=\bfseries},
    edcWarning/.style={colback=red!5, colframe=red!40!black, fonttitle=\bfseries},
    edcMechanism/.style={colback=teal!5, colframe=teal!50!black, fonttitle=\bfseries, title={Mechanistic Dimension Note (Canon)}}
}

\newcommand{\edcMechanismNote}[3]{%
\begin{tcolorbox}[edcMechanism]
\begin{itemize}[nosep,leftmargin=*]
    \item \textbf{5D cause (bulk):} #1
    \item \textbf{Brane-layer process:} #2
    \item \textbf{3D observation (output):} #3
\end{itemize}
\vspace{0.3em}
\footnotesize\textit{Ledger closure must hold: bulk + brane + 3D outputs conserve energy/quantum numbers.}
\end{tcolorbox}
}

% Chapter recap box
\newtcolorbox{chapterRecap}{
    colback=blue!3,
    colframe=blue!40!black,
    fonttitle=\bfseries,
    title={Chapter Recap},
    breakable
}

% ═══════════════════════════════════════════════════════════════════════════════
% TIKZ STYLES
% ═══════════════════════════════════════════════════════════════════════════════
\tikzset{
    edc box/.style={rectangle, draw, rounded corners=3pt, minimum width=2cm, minimum height=0.8cm, align=center, font=\small},
    edc arrow/.style={->, >=stealth, thick},
    gate box/.style={edc box, fill=orange!15, draw=orange!70!black, text=black},
    bulk region/.style={fill=gray!15, draw=none, rounded corners=0pt},
    brane region/.style={fill=blue!8, draw=none, rounded corners=0pt},
    observer region/.style={fill=green!8, draw=none, rounded corners=0pt},
    bulk box/.style={rectangle, rounded corners=4pt, draw=gray!60!black, thick, fill=gray!20, font=\footnotesize, align=center, minimum height=0.8cm},
    brane box/.style={rectangle, rounded corners=4pt, draw=blue!50!black, thick, fill=blue!15, font=\footnotesize, align=center, minimum height=0.8cm},
    process box/.style={rectangle, rounded corners=4pt, draw=orange!60!black, thick, fill=orange!15, font=\footnotesize, align=center, minimum height=0.8cm},
    output box/.style={rectangle, rounded corners=4pt, draw=green!50!black, thick, fill=green!15, font=\footnotesize, align=center, minimum height=0.8cm},
    section label/.style={font=\small\bfseries, text=black},
    bulk boundary/.style={draw=gray!60!black, thick, dashed},
    observer boundary/.style={draw=green!50!black, thick, dashed},
    brane boundary/.style={draw=blue!50!black, thick, dashed},
    particle/.style={circle, fill=green!40!black, minimum size=8pt, inner sep=0pt},
    neutrino/.style={circle, draw=green!40!black, fill=white, minimum size=8pt, inner sep=0pt},
    edc flow/.style={-{Stealth[length=5pt]}, thick, draw=black!70},
    edc dashed/.style={-{Stealth[length=4pt]}, thick, draw=gray!50, dashed},
}

% ═══════════════════════════════════════════════════════════════════════════════
% HYPERREF SETUP
% ═══════════════════════════════════════════════════════════════════════════════
\hypersetup{
  colorlinks=true,
  linkcolor=blue!60!black,
  citecolor=green!50!black,
  urlcolor=blue!70!black,
  pdftitle={EDC Book 2: Weak Sector (Canonical)},
  pdfauthor={Igor Grcman}
}

% ═══════════════════════════════════════════════════════════════════════════════
% SHARED MACROS
% ═══════════════════════════════════════════════════════════════════════════════
\input{_shared/meta/edc_meta_macros}
\input{_shared/meta/edc_stoplight_legend}

% ═══════════════════════════════════════════════════════════════════════════════
% DOCUMENT
% ═══════════════════════════════════════════════════════════════════════════════

\begin{document}

% ---------------------------------------------------------------------------
% FRONT MATTER
% ---------------------------------------------------------------------------
\frontmatter

% Title Page
\begin{titlepage}
\centering
\vspace*{2cm}
{\Huge\bfseries ELASTIC DIFFUSIVE COSMOLOGY}\\[1.5cm]
{\Large Part II: Weak Sector}\\[0.5cm]
{\large Geometric Origin of Weak Interactions}\\[2cm]
{\Large Igor Gr\v{c}man}\\[1cm]
{\large January 2026}\\[0.5cm]
{\small Canonical Edition}\\[2cm]
\vfill
{\large\itshape ``Weak interactions are not fundamental gauge vertices.\\
They are coarse-grained residues of thick-brane dynamics.''}
\end{titlepage}

% Copyright
\newpage
\thispagestyle{empty}
\vspace*{3cm}
\begin{center}
{\Large\bfseries Elastic Diffusive Cosmology}\\[0.3cm]
{\normalsize Quantum Mechanics and Gravity from a Unified 5D Membrane Action}\\[0.5cm]
{\large Part II: Weak Sector}\\[1.5cm]
\end{center}
\noindent\textcopyright{} 2026 Igor Gr\v{c}man\\[0.3cm]
All rights reserved under standard copyright.\\[0.5cm]
\noindent\textbf{DOI:} \href{https://doi.org/10.5281/zenodo.18328508}{10.5281/zenodo.18328508}\\[1cm]
\noindent Licensed under CC BY-NC-SA 4.0\\
\url{https://creativecommons.org/licenses/by-nc-sa/4.0/}\\[1.5cm]
\noindent\textbf{First edition:} January 2026\\
\vfill

% ---------------------------------------------------------------------------
% HOW TO READ THIS BOOK
% ---------------------------------------------------------------------------
\chapter*{How to Read This Book}
\addcontentsline{toc}{chapter}{How to Read This Book}

This book presents the EDC approach to weak interactions as a coherent narrative
in three parts. Each part serves a distinct purpose:

\begin{description}[leftmargin=2cm,labelwidth=1.8cm]
\item[Part I] \textbf{Physical Picture} --- Builds intuition through concrete examples.
    No equations beyond high-school algebra. If you read nothing else, read Part I.
\item[Part II] \textbf{Predictions} --- Derives electroweak observables from geometry.
    Requires comfort with basic quantum mechanics.
\item[Part III] \textbf{Derivation Machinery} --- Technical details for specialists.
    Contains the OPR (Open Problem Registry) chain with full derivations.
\end{description}

\section*{Epistemic Tags}

Every claim carries exactly one \textbf{evidence tag}:

\begin{center}
\begin{tabular}{@{}lp{10cm}@{}}
\toprule
\textbf{Tag} & \textbf{Meaning} \\
\midrule
\tagBL{}  & \textbf{Baseline} --- Accepted external fact (PDG/CODATA). Not derived. \\
\tagDer{} & \textbf{Derived} --- Follows explicitly from EDC postulates. \\
\tagDc{}  & \textbf{Derived conditional} --- Follows under stated assumptions. \\
\tagI{}   & \textbf{Identified} --- Structural mapping within EDC (not unique). \\
\tagP{}   & \textbf{Postulated} --- Hypothesis, not yet derived. \\
\tagM{}   & \textbf{Mathematics} --- Pure mathematical result (no physics input). \\
\bottomrule
\end{tabular}
\end{center}

\section*{The Projection Principle}

Every physical process has a \textbf{5D bulk+brane cause} whose observable
residue is a \textbf{3D shadow} on the observer boundary:

\begin{center}
\begin{tikzpicture}[scale=0.9]
  \node[bulk box, minimum width=3cm] (bulk) at (0,0) {5D Bulk Cause};
  \node[brane box, minimum width=3cm] (brane) at (4,0) {Brane Process};
  \node[output box, minimum width=3cm] (obs) at (8,0) {3D Shadow};
  \draw[edc flow] (bulk) -- (brane);
  \draw[edc flow] (brane) -- (obs);
\end{tikzpicture}
\end{center}

Standard Model terminology ($W^\pm$, $G_F$, V--A) appears as \textbf{3D observational
shorthand}. The SM describes \emph{what} observers measure; EDC explains \emph{why}.

\section*{Derivation Chain Map}

The book's logical dependencies form a directed graph:

\begin{center}
\begin{tikzpicture}[node distance=1.5cm, every node/.style={font=\small}]
  \node[edc box, fill=gray!20] (ch1) {Ch 1: Interface};
  \node[edc box, fill=gray!20, right=of ch1] (ch2) {Ch 2: Ontology};
  \node[edc box, fill=blue!10, below=of ch1] (ch3) {Ch 3--6: Cases};
  \node[edc box, fill=green!10, below=of ch2] (ch7) {Ch 7--8: EW Params};
  \node[edc box, fill=purple!10, right=of ch7] (ch9) {Ch 9--12: Mixing};
  \node[edc box, fill=orange!10, below=of ch3] (ch13) {Ch 13--16: OPR Chain};

  \draw[edc arrow] (ch1) -- (ch2);
  \draw[edc arrow] (ch2) -- (ch3);
  \draw[edc arrow] (ch3) -- (ch7);
  \draw[edc arrow] (ch7) -- (ch9);
  \draw[edc arrow] (ch3) -- (ch13);
  \draw[edc arrow] (ch13) -- (ch7);
\end{tikzpicture}
\end{center}

\textbf{Reading paths:}
\begin{itemize}[nosep]
  \item \textbf{Quick tour}: Ch 1--2, then Ch 7 summary boxes
  \item \textbf{Full narrative}: Ch 1--12 in order
  \item \textbf{Technical deep-dive}: Ch 13--16 after Part II
\end{itemize}

% ---------------------------------------------------------------------------
% TABLE OF CONTENTS
% ---------------------------------------------------------------------------
\tableofcontents

% ---------------------------------------------------------------------------
% MAIN MATTER
% ---------------------------------------------------------------------------
\mainmatter

% ###########################################################################
%                        PART I: PHYSICAL PICTURE
% ###########################################################################

\part{Physical Picture}
\label{part:physical_picture}

\begin{quote}
\textit{Part I builds intuition for the EDC approach to weak interactions.
No advanced mathematics required. Each chapter presents one particle decay
as a case study, showing how the unified pipeline applies.}
\end{quote}

% ===========================================================================
% CHAPTER 1: THE WEAK INTERFACE
% ===========================================================================
\chapter{The Weak Interface}
\label{ch:weak_interface}

\begin{quote}
\textit{This chapter establishes the physical mechanism for weak-sector processes.
We introduce the ``thick brane'' picture and explain why weak interactions appear
weak from a 5D perspective.}
\end{quote}

% --- Section 1.1: From 5D to Observable Physics ---
\section{From 5D Relaxation to Observable Physics}
\label{sec:5d_to_observable}

\input{sections/01_how_we_got_here}

% --- Section 1.2: The Unified Pipeline ---
\section{The Unified Pipeline: Absorption--Dissipation--Release}
% Label defined in sections/03_unified_pipeline.tex

\input{sections/03_unified_pipeline}

% --- Section 1.3: Geometry of the Brane Interface ---
\section{Geometry of the Brane Interface}
% Label defined in sections/02_geometry_interface.tex

\input{sections/02_geometry_interface}

% --- Chapter Recap ---
\begin{chapterRecap}
\textbf{What you learned:} Weak interactions are 3D shadows of 5D brane dynamics.
The unified pipeline (Absorption $\to$ Dissipation $\to$ Release) applies to all decays.

\textbf{What was derived:} Pipeline structure \tagDc{} from energy conservation.

\textbf{What was assumed:} Thick-brane geometry \tagP{}, proton stability \tagP{}.

\textbf{What remains open:} Quantitative mediator physics (Part III).

\textbf{Full proofs:} Derivation Library, ``Unified Pipeline'' section.
\end{chapterRecap}

% ===========================================================================
% CHAPTER 2: PARTICLE ONTOLOGY
% ===========================================================================
\chapter{Particle Ontology in EDC}
\label{ch:ontology}

\begin{quote}
\textit{EDC classifies particles into five ontological categories based on their
5D structure. This chapter defines each category and explains why the classification
matters for decay dynamics.}
\end{quote}

\input{sections/04_ontology}

% --- Section 2.2: Proton as Topological Anchor ---
\section{The Proton as Topological Anchor}
% Label defined in sections/04b_proton_anchor.tex

\input{sections/04b_proton_anchor}

% --- Section 2.3: Master Diagram ---
\section{The Master Diagram}
% Label defined in sections/04a_unified_master_figure.tex

\input{sections/04a_unified_master_figure}

% --- Chapter Recap ---
\begin{chapterRecap}
\textbf{What you learned:} Five particle categories (bulk-junction, brane-dominant,
brane-defect, edge-mode, composite) determine decay pathways.

\textbf{What was derived:} Classification scheme \tagDc{} from 5D structure.
Steiner theorem \tagM{} guarantees proton stability.

\textbf{What was assumed:} Y-junction topology \tagP{} for baryons.

\textbf{What remains open:} Quantitative mode profiles (Part III, OPR-21).

\textbf{Full proofs:} Derivation Library, ``Steiner Theorem'' and ``Z$_N$ Lemmas''.
\end{chapterRecap}

% ===========================================================================
% CHAPTER 3: NEUTRON DECAY
% ===========================================================================
\chapter{Case Study: Neutron Decay}
\label{ch:case_neutron}

\begin{quote}
\textit{The neutron provides our first test case: a bulk-junction particle
undergoing beta decay. This chapter traces the energy flow from junction
relaxation through brane charging to frozen projection.}
\end{quote}

\input{sections/05_case_neutron}

% Dual-route analysis (consolidated here)
\section{Route Analysis: Pump vs.\ Boundary-Dominant}
% Label defined in sections/05b_neutron_dual_route.tex

\input{sections/05b_neutron_dual_route}

% --- Chapter Recap ---
\begin{chapterRecap}
\textbf{What you learned:} Neutron decay is junction relaxation toward Steiner minimum.
The electron and antineutrino emerge from frozen projection of brane modes.

\textbf{What was derived:} Decay pathway \tagDc{} from pipeline + junction topology.

\textbf{What was assumed:} Neutron = displaced Y-junction \tagP{}.

\textbf{What remains open:} Numerical $\tau_n$ prediction (6\% achieved in Book 1).

\textbf{Full proofs:} Book 1, Chapter 3 (Peierls barrier calculation).
\end{chapterRecap}

% ===========================================================================
% CHAPTER 4: MUON DECAY
% ===========================================================================
\chapter{Case Study: Muon Decay}
\label{ch:case_muon}

\begin{quote}
\textit{The muon is a brane-dominant particle: its decay involves no bulk junction
dynamics. This simplifies the analysis and isolates the role of mode-profile
overlap in determining decay products.}
\end{quote}

\input{sections/06_case_muon}

% --- Chapter Recap ---
\begin{chapterRecap}
\textbf{What you learned:} Muon decay is internal brane-mode redistribution.
No bulk pumping needed---energy already stored in brane excitation.

\textbf{What was derived:} Decay mechanism \tagDc{} from brane-dominant ontology.

\textbf{What was assumed:} Mode-profile separation \tagP{}.

\textbf{What remains open:} Numerical $\tau_\mu$ from overlap integral.

\textbf{Full proofs:} Derivation Library, ``Projection-Reduction Lemma''.
\end{chapterRecap}

% ===========================================================================
% CHAPTER 5: TAU DECAY
% ===========================================================================
\chapter{Case Study: Tau Decay}
\label{ch:case_tau}

\begin{quote}
\textit{The tau has the same brane-dominant ontology as the muon but with higher
mode energy. This opens hadronic channels, providing a universality test for the
pipeline mechanism.}
\end{quote}

\input{sections/07_case_tau}

% --- Chapter Recap ---
\begin{chapterRecap}
\textbf{What you learned:} Tau decay confirms pipeline universality. Higher energy
opens hadronic channels via $\mathcal{P}_{\text{energy}}$ threshold.

\textbf{What was derived:} Same mechanism for all charged leptons \tagDc{}.

\textbf{What was assumed:} Mode index hierarchy \tagP{}.

\textbf{What remains open:} Quantitative branching ratios.

\textbf{Full proofs:} Same as muon (Chapter 4).
\end{chapterRecap}

% ===========================================================================
% CHAPTER 6: STABILITY AND EDGE MODES
% ===========================================================================
\chapter{Stability and Edge Modes}
\label{ch:stability_edge}

\begin{quote}
\textit{Three more particles complete the weak-sector catalog: the stable electron
(ground-mode brane defect), the pion (composite junction pair), and the neutrino
(edge mode at bulk-brane interface).}
\end{quote}

% --- Section 6.1: Electron Stability ---
\section{Electron: The Ground Mode}
% Label defined in sections/09_case_electron.tex

\input{sections/09_case_electron}

% --- Section 6.2: Pion Decay ---
\section{Pion: Composite Decay with Helicity Suppression}
% Label defined in sections/08_case_pion.tex

\input{sections/08_case_pion}

% --- Section 6.3: Neutrino Properties ---
\section{Neutrino: Edge-Mode Partner}
% Label defined in sections/10_case_neutrino.tex

\input{sections/10_case_neutrino}

% --- Chapter Recap ---
\begin{chapterRecap}
\textbf{What you learned:} Electron stability follows from being ground mode.
Pion helicity suppression follows from boundary conditions. Neutrinos are edge modes.

\textbf{What was derived:} Stability criterion \tagDc{}, helicity factor \tagDc{}.

\textbf{What was assumed:} Edge-mode localization \tagP{}.

\textbf{What remains open:} Neutrino mass mechanism.

\textbf{Full proofs:} Derivation Library, ``Chiral Localization'' section.
\end{chapterRecap}

% ###########################################################################
%                        PART II: PREDICTIONS
% ###########################################################################

\part{Predictions}
\label{part:predictions}

\begin{quote}
\textit{Part II derives electroweak observables from EDC geometry. We obtain
$\sin^2\theta_W$, $g^2$, $M_W$, and mixing angles without fitting to data.
Each result carries explicit epistemic tags showing derivation status.}
\end{quote}

% ===========================================================================
% CHAPTER 7: ELECTROWEAK PARAMETERS
% ===========================================================================
\chapter{Electroweak Parameters from Geometry}
\label{ch:electroweak}
\label{ch:z6_program}% Alias for legacy references

\begin{quote}
\textit{This chapter derives the core electroweak observables: Weinberg angle,
weak coupling, and W-boson mass. All emerge from the $\mathbb{Z}_6$ lattice
geometry established in Part I.}
\end{quote}

\input{CH3_electroweak_parameters}

% --- Chapter Recap ---
\begin{chapterRecap}
\textbf{What you learned:} $\sin^2\theta_W = 1/4$ (bare) from Z$_6$ geometry.
After RG running: 0.2314 vs.\ 0.2312 experiment (0.08\% error).

\textbf{What was derived:} $\sin^2\theta_W$ \tagDc{}, $g^2$ \tagDc{}, $M_W$ \tagDc{}.

\textbf{What was assumed:} Z$_6$ lattice structure \tagP{}.

\textbf{What remains open:} Full RG chain from first principles.

\textbf{Full proofs:} Derivation Library, ``Z$_N$ Discrete Averaging Lemma''.
\end{chapterRecap}

% ===========================================================================
% CHAPTER 8: THE FERMI CONSTANT (OVERVIEW)
% ===========================================================================
\chapter{The Fermi Constant: Overview}
\label{ch:gf_overview}

\begin{quote}
\textit{This chapter provides an overview of the $G_F$ derivation pathway.
The detailed machinery lives in Part III (Chapters 13--16). Here we explain
the structural argument: $G_F$ emerges from mode-overlap suppression.}
\end{quote}

\input{sections/11_gf_pathway}

% --- Chapter Recap ---
\begin{chapterRecap}
\textbf{What you learned:} $G_F$ is not a fundamental coupling but an overlap
suppression effect. The ``weakness'' of weak interactions is geometric.

\textbf{What was derived:} Structural formula $G_F \propto g_5^2 I_4 / m_\phi^2$ \tagDc{}.

\textbf{What was assumed:} Contact-interaction approximation \tagP{}.

\textbf{What remains open:} Numerical closure (Part III).

\textbf{Full proofs:} Chapter 15 (OPR-19, OPR-20, OPR-22).
\end{chapterRecap}

% ===========================================================================
% CHAPTER 9: THREE GENERATIONS
% ===========================================================================
\chapter{Why Exactly Three Generations?}
\label{ch:three_generations}

\begin{quote}
\textit{The Standard Model has three generations of quarks and leptons with no
explanation for this number. EDC connects it to the $\mathbb{Z}_3$ subgroup of
$\mathbb{Z}_6$ and the BVP bound-state spectrum.}
\end{quote}

\input{sections/05_three_generations}

% --- Chapter Recap ---
\begin{chapterRecap}
\textbf{What you learned:} Three generations arise from $\mathbb{Z}_3$ vortex modes
and the BVP bound-state count $N_{\text{bound}} = 3$ for appropriate $\mu$-window.

\textbf{What was derived:} Generation count constraint \tagDc{}.

\textbf{What was assumed:} BVP physical potential \tagP{}.

\textbf{What remains open:} Why $\mu$ falls in 3-generation window (Part III).

\textbf{Full proofs:} Chapter 14 (OPR-21 BVP closure).
\end{chapterRecap}

% ===========================================================================
% CHAPTER 10: NEUTRINO MIXING (PMNS)
% ===========================================================================
\chapter{Neutrino Mixing Angles}
\label{ch:pmns}
\label{ch:neutrinos_edge}% Alias for legacy references

\begin{quote}
\textit{The PMNS matrix describes neutrino flavor mixing. EDC derives the mixing
angles from $\mathbb{Z}_3 / \mathbb{Z}_6$ geometry. One angle achieves full derivation
status; two others are close.}
\end{quote}

\input{sections/06_neutrinos_edge_modes}

% --- Chapter Recap ---
\begin{chapterRecap}
\textbf{What you learned:} $\theta_{23}$ derived from Z$_6$ geometry (GREEN).
$\theta_{12}$, $\theta_{13}$ identified with small deviations (YELLOW).

\textbf{What was derived:} $\theta_{23} = 45°$ \tagDc{} (maximal mixing).

\textbf{What was identified:} $\theta_{12}$, $\theta_{13}$ \tagI{} with 8--9\% errors.

\textbf{What remains open:} Reactor angle from first principles.

\textbf{Full proofs:} Derivation Library, ``PMNS Attempts'' archive.
\end{chapterRecap}

% ===========================================================================
% CHAPTER 11: CKM AND CP VIOLATION
% ===========================================================================
\chapter{CKM Matrix and CP Violation}
\label{ch:ckm_cp}

\begin{quote}
\textit{Quark mixing and CP violation follow from the same $\mathbb{Z}_3$ geometry
that governs leptons. The CP phase emerges from $\mathbb{Z}_2$ selection within
the $\mathbb{Z}_6$ lattice.}
\end{quote}

\input{sections/07_ckm_cp}

% --- Chapter Recap ---
\begin{chapterRecap}
\textbf{What you learned:} CP phase $\delta \approx 60°$ from Z$_2$ selection (5° from PDG).
Pure Z$_3$ gives $J_{CP} = 0$ (Phase Cancellation Theorem).

\textbf{What was derived:} Phase Cancellation Theorem \tagDc{}.

\textbf{What was identified:} $\delta = 60°$ \tagI{} from Z$_2$ selection.

\textbf{What remains open:} Jarlskog invariant from geometry.

\textbf{Full proofs:} Derivation Library, ``CP Phase Attempts'' archive.
\end{chapterRecap}

% ===========================================================================
% CHAPTER 12: V-A STRUCTURE
% ===========================================================================
\chapter{V--A Structure from Chiral Localization}
\label{ch:va_structure}

\begin{quote}
\textit{The V--A structure of weak currents is not postulated but emerges from
5D chiral localization. Left-handed and right-handed modes have different profiles,
leading to chiral selection at the boundary.}
\end{quote}

\input{sections/09_va_structure}

% --- Chapter Recap ---
\begin{chapterRecap}
\textbf{What you learned:} V--A emerges from asymmetric L/R mode profiles.
Boundary conditions project onto left-handed states.

\textbf{What was derived:} Chiral selection mechanism \tagDc{}.

\textbf{What was assumed:} Domain-wall mass profile \tagP{}.

\textbf{What remains open:} Right-handed current suppression factor.

\textbf{Full proofs:} Derivation Library, ``Chiral Localization'' section.
\end{chapterRecap}

% ###########################################################################
%                        PART III: DERIVATION MACHINERY
% ###########################################################################

\part{Derivation Machinery}
\label{part:machinery}

\begin{quote}
\textit{Part III contains the technical OPR (Open Problem Registry) derivation
chain. This is specialist material. The key dependencies are:
Scale Taxonomy (OPR-04) $\to$ BVP Framework (OPR-21) $\to$ Coupling Chain
(OPR-19/20/22) $\to$ $G_F$ Closure.}
\end{quote}

% ===========================================================================
% CHAPTER 13: SCALE TAXONOMY AND ANCHORS
% ===========================================================================
\chapter{Scale Taxonomy and Anchors}
\label{ch:scale_taxonomy}

\begin{quote}
\textit{This chapter defines all length/energy scales precisely and derives
the $\sigma \to M_0$ anchor. These definitions are prerequisites for the
BVP framework in Chapter 14.}
\end{quote}

% OPR-04: Scale Taxonomy
\section{OPR-04: Scale Taxonomy}
\label{sec:opr04}

\input{sections/ch16_opr04_delta_derivation}

% OPR-01: Sigma Anchor
\section{OPR-01: The $\sigma \to M_0$ Anchor}
\label{sec:opr01}

\input{sections/ch15_opr01_sigma_anchor_derivation}

% --- Chapter Recap ---
\begin{chapterRecap}
\textbf{What you learned:} Four scales ($\Delta$, $\delta$, $\ell$, $R_\xi$) have
distinct physical meanings. The $\sigma \to M_0$ anchor reduces free parameters.

\textbf{What was derived:} Scale taxonomy \tagDer{}, $M_0$ formula \tagDer{}.

\textbf{What was assumed:} Assumption labels (A1)--(A3) for scale identifications.

\textbf{What remains open:} Which assumption (A1/A2/A3) nature chooses.

\textbf{Full proofs:} Derivation Library, ``OPR Registry Summary''.
\end{chapterRecap}

% ===========================================================================
% CHAPTER 14: BVP FRAMEWORK
% ===========================================================================
\chapter{BVP Framework}
\label{ch:bvp_framework}
\label{ch:bvp_master_key}% Alias for legacy references

\begin{quote}
\textit{The Boundary Value Problem (BVP) is the mathematical core of EDC weak-sector
predictions. This chapter specifies the problem and derives the Robin boundary
conditions from action variation.}
\end{quote}

% Infrastructure
\section{BVP Specification}
\label{sec:bvp_spec}

\input{sections/ch12_bvp_workpackage}

% OPR-21: BVP Closure
\section{OPR-21: BVP Closure}
\label{sec:opr21}

\input{sections/ch14_opr21_closure_derivation}

% Electroweak Bridge
\section{The Electroweak Bridge}
\label{sec:electroweak_bridge}

\input{sections/ch10_electroweak_bridge}

% --- Chapter Recap ---
\begin{chapterRecap}
\textbf{What you learned:} BVP determines mode profiles, eigenvalues, and overlap
integrals. Robin BC emerges from action variation. Three bound states achieved
for $\mu \in [\mu_3^-, \mu_3^+]$.

\textbf{What was derived:} Robin BC form \tagDer{}, eigenvalue structure \tagDer{}.

\textbf{What was assumed:} Domain-wall potential shape \tagP{}.

\textbf{What remains open:} $\alpha$ numerical value (OPR-20b).

\textbf{Full proofs:} Derivation Library, ``Frozen Brane BC Complete Analysis''.
\end{chapterRecap}

% ===========================================================================
% CHAPTER 15: COUPLING CHAIN
% ===========================================================================
\chapter{The Coupling Chain: $g_5 \to G_F$}
\label{ch:coupling_chain}
\label{ch:gf_derivation}% Alias for legacy references

\begin{quote}
\textit{This chapter completes the derivation of $G_F$ by assembling the coupling
chain: 5D coupling $g_5$ $\to$ mediator mass $m_\phi$ $\to$ effective coupling
$G_{\text{eff}}$.}
\end{quote}

% OPR-19: g5 from Action
\section{OPR-19: The 5D Coupling $g_5$}
\label{sec:opr19}

\input{sections/ch17_opr19_g5_from_action}

% OPR-20: Mediator Mass
\section{OPR-20: Mediator Mass from Eigenvalue}
\label{sec:opr20}

\input{sections/ch18_opr20_mediator_mass_from_eigenvalue}

% OPR-22: G_eff from Exchange
\section{OPR-22: $G_{\text{eff}}$ from Exchange}
\label{sec:opr22}

\input{sections/ch19_opr22_geff_from_exchange}

% --- Chapter Recap ---
\begin{chapterRecap}
\textbf{What you learned:} $G_F = g_5^2 |f_1(0)|^2 / (2 x_1^2)$ from mediator exchange.
The ``weakness'' of weak interactions is overlap suppression.

\textbf{What was derived:} $g_5$ normalization \tagDer{}, exchange formula \tagDer{}.

\textbf{What was assumed:} Contact approximation \tagP{}.

\textbf{What remains open:} OPR-20b ($\alpha$ provenance) has YELLOW status.

\textbf{Full proofs:} Derivation Library, ``$G_F$ Non-Circular Chain''.
\end{chapterRecap}

% ===========================================================================
% CHAPTER 16: CLOSURE STATUS
% ===========================================================================
\chapter{Closure Status and Open Problems}
\label{ch:closure_status}

\begin{quote}
\textit{This chapter consolidates the epistemic status of all Part III derivations
and catalogs remaining open problems.}
\end{quote}

\input{sections/ch20_epistemic_summary_closure_status}

% Consolidated gate registry and epistemic map (provides sec:gate_registry)
\input{sections/12_epistemic_map}

% --- Chapter Recap ---
\begin{chapterRecap}
\textbf{What you learned:} Most electroweak predictions achieve GREEN or YELLOW status.
$G_F$ chain is structurally complete; numerical closure awaits BVP inputs.

\textbf{Closure status:}
\begin{itemize}[nosep]
  \item $\sin^2\theta_W$: GREEN \tagDc{} (0.08\% error)
  \item $M_W$: GREEN \tagDc{} (0.2\% error)
  \item $\theta_{23}$: GREEN \tagDc{}
  \item $G_F$: YELLOW \tagDc{} (structure complete, numerics pending)
  \item OPR-20b ($\alpha$): YELLOW/RED-C (route A blocked)
\end{itemize}

\textbf{What remains open:} See OPR Registry in Derivation Library.
\end{chapterRecap}

% ###########################################################################
%                           EPILOGUE
% ###########################################################################

\part*{Epilogue}
\addcontentsline{toc}{part}{Epilogue}

% ===========================================================================
% CHAPTER 17: NUCLEAR APPLICATIONS PREVIEW
% ===========================================================================
\chapter{Nuclear Applications Preview}
\label{ch:nuclear_preview}

\begin{quote}
\textit{This chapter previews nuclear applications developed in Book 3:
the MN topology, the frustration-corrected Geiger--Nuttall law, and the
44.7\% discrepancy resolution for alpha-decay lifetime predictions.}
\end{quote}

\input{sections/XX_teaser_book3_nuclear_mn_gn}

% ---------------------------------------------------------------------------
% BACK MATTER
% ---------------------------------------------------------------------------
\backmatter

% Derivation Library
\input{appendices/APPENDIX_DERIVATION_LIBRARY}

% EDC vs SM Comparison (moved to appendix)
\chapter*{Appendix: EDC vs Standard Model}
\addcontentsline{toc}{chapter}{Appendix: EDC vs Standard Model}

\begin{center}
\begin{longtable}{p{4cm}p{5cm}p{5cm}}
\toprule
\textbf{Aspect} & \textbf{Standard Model} & \textbf{EDC} \\
\midrule
\endfirsthead
\toprule
\textbf{Aspect} & \textbf{Standard Model} & \textbf{EDC} \\
\midrule
\endhead
\bottomrule
\endfoot
\bottomrule
\endlastfoot

Spacetime & 4D Minkowski & 5D bulk + 4D brane \\
Gauge structure & Postulated & Emerges from geometry \\
$\sin^2\theta_W$ & Input (0.2312) & Derived: 0.2314 (0.08\%) \\
$G_F$ & Input & Derived from overlap \\
Why V--A? & Postulated & Chiral localization \\
Why 3 generations? & No explanation & Z$_3$ + BVP spectrum \\
\end{longtable}
\end{center}

% Notation
\chapter*{Appendix: Notation and Symbols}
\addcontentsline{toc}{chapter}{Appendix: Notation}

\begin{center}
\begin{tabular}{lll}
\toprule
\textbf{Symbol} & \textbf{Meaning} & \textbf{Typical Value} \\
\midrule
$\sigma$ & Membrane tension & $5.86 \times 10^{6}$ MeV/fm$^2$ \\
$r_e$ & Lattice spacing & 2.82 fm \\
$\delta$ & Boundary-layer scale & $\sim r_e$ \\
$\Delta$ & Kink width & $\sim 10^{-3}$ fm \\
$\ell$ & Domain support & $\sim 1$ fm \\
$R_\xi$ & Diffusion scale & $\hbar c / M_Z$ \\
\bottomrule
\end{tabular}
\end{center}

% Bibliography
\printbibliography[heading=bibintoc,title={References}]

\end{document}
